\chapter{Thảo luận và giới hạn}

Kết quả mạnh nhất của luận văn là độ trễ và độ tin cậy của việc đo trên cụm thật. Giới hạn lớn nhất là recall đo được chỉ là cận trên cho tấn công đi vào qua exec, và detector chưa phân biệt được ác ý với sự mới lạ.

\section{Detector phát hiện sự mới lạ, không phải ác ý}\label{sec:limits-novelty}

Mọi tấn công trong campaign đều là một process mới chạy qua exec, và Pulse bắt “có process mới” rất mạnh. Hai quan sát cho thấy điều này:

\begin{itemize}
\tightlist
\item
  binary vô hại bị gắn cờ 80,7--95,1\%;
\item
  57,9\% cảnh báo trên normal trùng lúc có exec.
\end{itemize}

Vì vậy recall gần 100\% có thể bị thổi phồng một phần bởi chính dấu hiệu exec. Đây là dạng tương quan giả mà \citet{arp2022} cảnh báo. Đối chứng binary vô hại đo được mức thổi phồng, nhưng không cho biết recall khi tấn công không sinh process mới.

Tripwire tách được một lớp: tấn công buộc phải dùng syscall lạ. Hai kịch bản \texttt{mprotect} và \texttt{seccomp} thì chỉ còn tín hiệu tần suất. Muốn tiến xa hơn cần thông tin ngoài số đếm, ví dụ tham số syscall hay quan hệ cha--con của process (Chương 8).

\section{Gánh cảnh báo}

Gánh 1,50--1,55 đợt/replica-giờ còn gấp khoảng 15 lần ngân sách 0,1. Với khoảng 80 container, con số này tương đương cỡ 2.900 đợt cảnh báo mỗi ngày (ước tính 1,5 × 80 × 24).

Precision 99,4\% đo trên campaign cân bằng 1:1. Ở tỉ lệ tấn công thật, precision sẽ thấp hơn rất nhiều. Vì vậy luận văn báo gánh cảnh báo như một giới hạn, không coi precision là bằng chứng dùng được trong vận hành.

Phần lớn cảnh báo còn lại là benign trigger có cấu trúc, nên có thể xử lý ở tầng sau:

\begin{itemize}
\tightlist
\item
  probe exec của Kubernetes;
\item
  một replica Kafka lệch so với hai broker còn lại.
\end{itemize}

Lọc probe đã giảm 23\% gánh, nhưng phải chờ RCA nên chậm. Calibration riêng theo replica cho Kafka chưa được thử.

Về độ lớn: muốn về ngân sách 0,1 thì phải giảm thêm khoảng 93,5\%. Mức giảm báo nhầm tốt nhất đã công bố bằng cách xếp hạng cảnh báo theo ngữ cảnh, khoảng 84--86\% (\citealp{hassan2019}), vẫn để lại 2--2,5 lần ngân sách. Khoảng cách này là hệ quả của việc đo sự mới lạ, không phải của việc chỉnh tham số. Nghiên cứu về SOC cũng cho thấy phần lớn cảnh báo bị coi là ``sai'' thực ra là benign trigger: hệ thống thấy đúng sự kiện, chỉ là sự kiện vô hại (\citealp{alahmadi2022}). Gom các đợt cảnh báo thành sự cố và tương quan với sự kiện của Kubernetes là đòn bẩy còn lại cho vận hành, nhưng không thay đổi tỉ lệ binary vô hại bị gắn cờ.

\section{Tính thực tế của tấn công}

\begin{itemize}
\tightlist
\item
  \textbf{Giai đoạn:} fixture tự viết mô phỏng giai đoạn sau xâm nhập. Mọi thao tác cố ý thất bại hoặc chỉ ảnh hưởng chính process, kéo dài 15 s, với tốc độ 1--24 lần/s.
\item
  \textbf{Độ phủ ATT\&CK for Containers:} chưa có initial access qua ứng dụng, persistence, credential access, discovery, lateral movement hay impact (ví dụ đào coin).
\item
  \textbf{Tấn công đi qua process của dịch vụ} (SQL injection, command injection) chưa được đo trên cụm.
\item
  \textbf{Exploit thật:} đánh giá trên LID-DS (mục 6.12) bù một phần với 15 kịch bản lỗ hổng thật, nhưng đó là đánh giá offline trên trace ghi sẵn, không phải trên cụm.
\end{itemize}

Về quy mô, campaign của luận văn không nhỏ so với các đánh giá đã công bố: một so sánh gần đây giữa Falco, Tetragon và Tracee dùng 3 tấn công, mỗi tấn công lặp 20 lần, và báo tỉ lệ báo nhầm 0\% mà không mô tả tải vô hại (\citealp{syairozi2025}); ReplicaWatcher dùng 13 kịch bản CVE/CWE trên hai ứng dụng (\citealp{elkhairi2024}). Khoảng trống thật của luận văn không nằm ở số kịch bản mà ở loại kịch bản: chưa có exploit chạy trong chính process của dịch vụ trên cụm, đúng loại mà LID-DS cho thấy Pulse yếu nhất.

\section{Kẻ tấn công thích nghi}

\begin{itemize}
\tightlist
\item
  \textbf{Dùng syscall có sẵn:} kẻ tấn công chỉ dùng syscall đã có trong hồ sơ normal thì tránh được tripwire, và có thể “pha loãng” tín hiệu tần suất (\citealp{wagner2002}). Theo khung của \citet{tan2002}, một detector chuỗi chỉ thấy một chuỗi lạ khi cửa sổ của nó dài ít nhất bằng chuỗi đó. Tripwire là trường hợp cửa sổ dài 1: nó chỉ thấy ID syscall lạ, nên mọi tấn công ghép từ syscall quen đều lọt. Né tránh kiểu pha loãng đã được làm tự động với độ thành công 100\% trên năm hệ phát hiện dựa trên provenance (\citealp{goyal2023}).
\item
  \textbf{Chạy chậm:} ở 1 lần/s, recall còn 79,0\% ở lượt cuối, so với 93,1\% ở lượt \texttt{pulse-\allowbreak{}stable-\allowbreak{}confirm} trước khi có chiết khấu theo replica. Hai số đến từ hai lượt khác nhau, nên chiết khấu là nguyên nhân khả dĩ, chưa được tách riêng trên cùng dữ liệu.
\item
  \textbf{Tấn công mọi replica cùng lúc:} chiết khấu chỉ áp vào nhánh ML. Tripwire vẫn báo nếu tấn công dùng ID syscall lạ. Với quyền \texttt{kubectl\ exec}, chạy cùng một lệnh trên mọi replica là việc rẻ, khác giả định của ReplicaWatcher rằng kẻ tấn công khó làm điều đó (\citealp{elkhairi2024}).
\item
  \textbf{Đi vòng tầng syscall:} thao tác gửi qua \texttt{io\_\allowbreak{}uring} không đi qua \texttt{sys\_\allowbreak{}enter}, nên công cụ chỉ móc syscall không thấy chúng (\citealp{armo2025}). Pulse đếm mọi ID syscall, nên lời gọi \texttt{io\_\allowbreak{}uring\_\allowbreak{}setup} sẽ bị tripwire bắt nếu workload chưa từng dùng nó; các thao tác sau đó thì Pulse không thấy. Điều này chưa được kiểm trên cụm.
\end{itemize}

Luận văn chưa làm thí nghiệm né tránh có chủ đích. Một thí nghiệm tối thiểu cần chia kẻ tấn công theo mức hiểu biết về detector (không biết, biết tập syscall, biết cả chính sách chiết khấu) và đo cả tỉ lệ đạt mục tiêu của tấn công, không chỉ tỉ lệ bị bắt.

\section{Calibration và giả thiết hoán đổi}\label{sec:limits-calibration}

P-value conformal chỉ hợp lệ khi normal lúc chạy giống normal lúc hiệu chỉnh. Trên cụm thật giả thiết này bị phá:

\begin{itemize}
\tightlist
\item
  calibration cũ dần theo ngày;
\item
  một sự cố thật (replica CloudNativePG hỏng sau failover) tạo ra cảnh báo liên tục.
\end{itemize}

Vì vậy ngân sách 5\% chỉ là ngân sách thiết kế, không phải bảo đảm.

Calibration tự làm mới đã chạy smoke đạt và đang được xác nhận live (mục~\ref{sec:live-recal}). Khi chạy, nó sẽ loại khoảng campaign khỏi calibration như bản offline. Một hệ thống vận hành thật không biết lúc nào bị tấn công, nên câu hỏi là điều gì xảy ra khi cửa sổ tấn công lọt vào calibration:

\begin{itemize}
\tightlist
\item
  nếu điểm của cửa sổ tấn công cao hơn điểm normal, p-value conformal trở nên bảo thủ: tỉ lệ báo nhầm vẫn giữ, cái mất là recall (\citealp{bashari2025}). Với ngân sách α = 0,00167 mỗi cửa sổ, chỉ cần một tấn công khoảng 10 phút trong 12 giờ calibration là đủ thay phần đuôi của phân phối (suy luận, chưa đo);
\item
  ngược lại, tự bỏ các cửa sổ đã báo khỏi calibration mà không có nhãn làm tỉ lệ báo nhầm phình 1,7--3,6 lần mức danh nghĩa trong thí nghiệm của cùng nhóm tác giả; chỉ bỏ các cửa sổ đã được xác nhận là tấn công thì mức phình bị chặn ở khoảng $1/(n_0+1)$;
\item
  các bộ điều khiển tần suất báo như ACI (\citealp{gibbs2021}) hay conformal PID (\citealp{angelopoulos2023}) giữ tỉ lệ báo trung bình đúng ngân sách với mọi chuỗi điểm, nhưng vì thế sẽ nâng dần ngưỡng trong một tấn công kéo dài và che nó đi.
\end{itemize}

\section{Phạm vi của bằng chứng}

\begin{itemize}
\tightlist
\item
  \textbf{Một cụm, một ứng dụng, một kernel:} Ubuntu 24.04, kernel 6.8, x86-64, chạy trên VM. Overhead và độ trễ có thể khác trên phần cứng hay kernel khác. Luận văn không kiểm trên kernel hoặc distro khác.
\item
  \textbf{Nhãn:} là nhãn kỳ vọng theo thiết kế trial. Cảnh báo trên normal được review, nhưng không có người review độc lập.
\item
  \textbf{Kết quả development:} một số kết quả (nhóm feature, baseline) là development và được ghi rõ như vậy. Chỉ kết quả blind được dùng để kết luận cuối.
\item
  \textbf{Đoạn normal ngắn:} mỗi lượt live chỉ có 24 giờ, và các chế độ tải biến động tập trung trong 90 phút đầu.
\end{itemize}

\section{Truy nguyên nguyên nhân}

Cây bằng chứng cho thấy phụ thuộc theo định danh và thời gian, yếu hơn phụ thuộc theo luồng thông tin mà các hệ provenance dùng, và không chứng minh nhân quả. Khi hai process mới chạy gần nhau và bằng điểm, luật cũ phá hoà theo tên binary, nên kết quả phụ thuộc cách đặt tên; luật phá hoà theo thời điểm khởi chạy sửa được artifact này trên campaign, nhưng sẽ xếp sai khi tấn công khởi chạy trước một process mới vô hại. Khoảng 11\% tấn công không có case riêng, và nguyên nhân chính còn là giả thuyết chưa kiểm chứng: sự kiện exec của tấn công tới RCA sau cảnh báo. Tetragon tự ghi trong mã nguồn rằng sự kiện có thể tới trước sự kiện exec của chính process, và giữ lại để thử lại mỗi 2 s, tối đa 15 lần (\citealp{tetragon_eventcache}), nên một sự kiện có thể trễ tới khoảng 30 s. Đây là cơ chế ứng viên, chưa đo trên cụm.
